\section{Sustav preporuke sljedećeg pitanja}
U sustavima dijaloga, problem preporuke sljedećeg pitanja (ili rangiranja sljedećeg iskaza) definira se kao zadatak odabira najprikladnijeg sljedećeg iskaza $u_{n+1}$ (odgovora) iz unaprijed definiranog skupa od $m$ kandidatskih iskaza $\{r_1, r_2,..., r_m\}$, s obzirom na kontekst razgovora $C = \{u_1, u_2,..., u_n\}$ (slijed od $n$ iskaza). Cilj je rangirati ove kandidatske odgovore na temelju njihove prikladnosti za kontekst razgovora, pri čemu se bira onaj s najvišim rezultatom.

Ovaj zadatak predstavlja nekoliko značajnih izazova:
\begin{itemize}
\item \textbf{Implicitna zajednička informacija:} Često postoji implicitna zajednička informacija između konteksta i odgovora koja nije eksplicitno navedena.
\item \textbf{Identificiranje važnih informacija:} Identificiranje važnih informacija (riječi, fraza i rečenica) unutar konteksta i njihovo učinkovito podudaranje s odgovarajućim informacijama u odgovoru je složeno.
\item \textbf{Modeliranje odnosa između višestrukih iskaza:} Modeliranje odnosa između višestrukih iskaza koji čine povijest razgovora također je značajan izazov. Iako neki pristupi pojednostavljuju razgovor na samo posljednji iskaz, uzimanje u obzir **cjelokupnog konteksta** je ključno za bolju izvedbu.
\item \textbf{Ograničeno razumijevanje jezika:} Chatbotovi se suočavaju s poteškoćama u razumijevanju nijansi poput idioma, sarkazma ili igre riječi, što rezultira pogrešnim tumačenjima i irelevantnim odgovorima. Čak i napredni NLP modeli mogu pogrešno protumačiti polisemne pojmove ili regionalne dijalekte.
\item \textbf{Nemogućnost učenja iz iskustva:} Tradicionalni chatbotovi ne mogu prilagoditi odgovore na temelju povijesti razgovora niti poboljšati performanse tijekom vremena.
\item \textbf{Nedostatak emocionalne inteligencije:} Chatbotovi ne mogu prepoznati korisnikovo emocionalno stanje, što može dovesti do neprikladnih ili neosjetljivih odgovora.
\item \textbf{Problem hladnog starta:} Nedostatak visokokvalitetnih povijesnih podataka za nove domene dovodi do loše klasifikacije namjera i irelevantnih putanja ispitivanja.
\item \textbf{Pristranost uzorka u obrascima odgovora:} Chatbotovi mogu angažirati samo digitalno pismene korisnike, iskrivljujući podatke i uvodeći demografsku/bihevioralnu pristranost.
\item \textbf{Slaba dugoročna memorija u LLM-ovima:} LLM-ovi koji pokreću chatbotove često pokazuju slabo dugoročno pamćenje, zaboravljajući ranije unose ili proturječeći sami sebi tijekom razgovora.
\item \textbf{Rukovanje neviđenim scenarijima:} Chatbotovi se često susreću s nepredviđenim scenarijima, što može dovesti do "halucinacija" - generiranja netočnih informacija s visokim stupnjem samopouzdanja.
\end{itemize}


\subsection{Temeljni Principi i metodologije}
Temeljni principi rangiranja sljedećeg pitanja uključuju sofisticirane metode reprezentacije teksta, arhitekture kodiranja i matematičke operacije za izračun sličnosti.

\textbf{Reprezentacija Teksta (Vektorski Ugradci):} Ulazni tekst, uključujući kontekst razgovora i kandidatski odgovor, prvo se pretvara u vektorske ugradke. To su numeričke vektorske reprezentacije riječi, gdje riječi sa sličnim značenjima imaju slične vektorske reprezentacije. Ovi ugradci se često inicijaliziraju unaprijed obučenim modelima poput GloVe-a ili se generiraju modelima kao što su Word2Vec, BERT, Sentence-BERT i Universal Sentence Encoder.

\textbf{Arhitekture Kodiranja (npr. LSTM Dual Enkoderi, Pažljivi Dual Enkoderi):}
\begin{itemize}
\item \textbf{Dual Enkoder Model:} Sastoji se od dva odvojena enkodera (npr. Long Short-Term Memory, LSTM, mreže) koji neovisno izvlače značajke za dijalog kontekst i kandidatski odgovor. Ovi enkoderi transformiraju ulazni tekst varijabilne veličine (sekvencu vektorskih ugradaka riječi) u vektore fiksne duljine, označene kao $C'$ za kontekst i $R'$ za odgovor. Skriveni sloj LSTM enkodera ažurira se svaki put kada se obradi vektorski ugradak riječi.

\item \textbf{Pažljivi Dual Enkoder Model:} Predstavlja napredak u odnosu na standardne dual enkodere, uključujući mehanizam pažnje povrh izvučenih značajki na razini riječi iz oba enkodera. Ovaj mehanizam uči vezu između dijalog konteksta i odgovora na razini riječi. Matrica sličnosti $S$ (npr. koristeći kosinusnu sličnost) definira se na kodiranim značajkama $h_x$ (kontekst) i $h_y$ (odgovor) za mjerenje parnih odnosa riječi.[19] Težine pažnje se izračunavaju (npr. korištenjem softmax funkcije preko eksponencijala unosa matrice sličnosti) i zatim ponderiraju izvorne kodirane značajke prema njihovoj važnosti.
\end{itemize}



\subsection{Matematička osnova izračuna sličnosti}
\subsubsection{Vektorski produkt (u nekim dual enkoderima)}
U nekim poboljšanim arhitekturama LSTM dual enkodera, sličnost između kodiranog kontekstualnog vektora $C'$ i kodiranog vektora odgovora $R'$ izračunava se kao križni produkt $P$. Matematički, ako je $C' = [c_1, c_2,..., c_h]$ i $R' = [r_1, r_2,..., r_h]$ (gdje je $h$ veličina skrivenog sloja), tada je $$P = C' \times R' = [c_1 \cdot r_1, c_2 \cdot r_2,..., c_h \cdot r_h]$$ (element-po-element množenje, često se naziva Hadamardovim produktom u ovom kontekstu). Cilj ovog produkta je uhvatiti zajedničke semantičke i sintaktičke značajke.

\subsubsection{Skalarni produkt (u attentive dual encoderima)}
Konačni rezultat predviđanja za par dijalog kontekst-odgovor u attentive dual enkoderima izračunava se kao skalarni produkt pažljivo ponderiranih značajki konteksta i odgovora: $$score(x, y) = dot((a_x)|h_x, (a_y)|h_y).$$

\subsubsection{Generiranje vjerojatnosnog rezultata} Vektorski rezultat sličnosti (npr. $P$ iz križnog produkta) ili rezultat (iz skalarnog produkta) prosljeđuje se u potpuno povezani sloj. Na izlaz ovog sloja primjenjuje se sigmoidna aktivacijska funkcija, $\sigma(x) = \frac{1}{1 + e^{-x}}$, kako bi se dobio vjerojatnosni rezultat između 0 i 1. Ovaj rezultat ukazuje na vjerojatnost da je kandidatski iskaz ispravan sljedeći iskaz za zadani kontekst i koristi se za rangiranje.

\subsection{Poboljšanje interpretativnosti (attentive dual enkoderi)}
\textbf{Problem:} Standardni mehanizmi pažnje mogu istaknuti nevažne riječi, zamagljujući interpretativnost i potencijalno narušavajući performanse. Enkoderi također mogu izgubiti fine-zrnate informacije na razini riječi.

\begin{itemize}
\item Rješenje 1: \textbf{Regularizacija bez pažnje (minimizacija uzajamne informacije):} Predlaže se novi gubitak regularizacije koji minimizira uzajamnu informaciju između nevažnih riječi konteksta ($h_x$) i značajki odgovora ($h_y$). To prisiljava model da de-naglaava nevane rijei dok naglaava vane.

\item Rješenje 2: \textbf{Rezidualni sloj s sirovim vektorskim ugradcima riječi:} Dodaje se rezidualni sloj koji izravno povezuje sirove vektorske ugradke riječi ($e_x$) s konačnom kodiranom značajkom konteksta, čime se čuva interpretativnost na razini riječi.
\end{itemize}


\subsection{Ukupni Cilj Treninga (attentive dual enkoderi)} Ukupni cilj treninga je min-max igra koja kombinira gubitak dohvaćanja ($L_{ret}$) za maksimiziranje rezultata ispravnih parova dijalog kontekst-odgovor i regularizacijski član uzajamne informacije ($L_{reg}$) za isticanje korisnih informacija: $$\min_{\theta_x,\theta_y} \max_{\theta_D} L_{ret}(\theta_x, \theta_y) + \beta L_{reg}(\theta_x, \theta_y, \theta_D).$$


Sustav preporuke može koristiti istu tehniku semantičke sličnosti kao i Semantic Router. 
Algoritam radi kako slijedi:
\begin{itemize}
    \item Nakon što korisnik postavi upit i dobije odgovor, sustav generira vektorsku reprezentaciju upita ($v_q$).
    \item Uspoređuje ($v_q$) s vektorima svih pitanja ($v_{p_i}$) u bazi podataka.
    \item Odabire $k$ pitanja (npr. $k = 3$) s najvećom kosinusnom sličnošću, isključujući trenutno pitanje.
    \item Prikazuje preporučena pitanja korisniku u sučelju.
\end{itemize}

Ako korisnik pita "Koji su uvjeti za upis?", sustav može predložiti pitanja poput:
\begin{itemize}
\item   "Kako se prijaviti za stipendiju?"
\item   "Kada su rokovi za upis?"
\item   "Koje dokumente treba dostaviti za upis?"
\end{itemize}



\subsection{Buduće nadogradnje}

\begin{itemize}
    \item \textbf{Personalizacija:} Sustav može pratiti povijest korisnikovih upita i prilagođavati preporuke na temelju njihovih interesa.
    \item \textbf{Kontekstualna analiza:} Umjesto samo semantičke sličnosti, sustav može koristiti dodatne informacije, poput kategorija pitanja (npr. upis, stipendije, ispiti).
    \item \textbf{Interaktivni prijedlozi:} Omogućiti korisnicima ocjenjivanje preporuka kako bi se poboljšao algoritam.
\end{itemize}
